\documentclass{article}
\usepackage{amsmath,amssymb,amsthm}
\usepackage{graphicx}
\usepackage{algorithm}
\usepackage{algorithmic}
\usepackage{mathtools}
\usepackage{hyperref}
\usepackage{xcolor}

\title{Funding Rate Arbitrage: Mathematical Framework \\ \Large{(Updated with Critique Integrations)}}
\author{CyberDeltaEngine Project}
\date{\today}

\newtheorem{theorem}{Theorem}
\newtheorem{lemma}[theorem]{Lemma}
\newtheorem{proposition}[theorem]{Proposition}
\newtheorem{corollary}[theorem]{Corollary}
\newtheorem{definition}{Definition}
\newtheorem{assumption}{Assumption}

\begin{document}

\maketitle

\begin{abstract}
This document formalizes the mathematical framework for funding rate arbitrage strategies across perpetual futures markets, based on research from He et al. (2024), Alexander et al. (2024), and Fanelli et al. (2024). We derive no-arbitrage bounds, analyze funding mechanics, and develop optimal strategies. Special attention is given to Hyperliquid (hourly funding, Hyperps). This revision incorporates advanced concepts based on critique: refined no-arbitrage bounds including execution costs, stochastic and adaptive modeling for funding rates and volatility, refined Kelly criterion (\textbf{including transfer cost considerations and constraints}), explicit modeling of basis risk and optimal dynamic hedging, and enhanced risk management considerations (e.g., CVaR). \textbf{It should be read in conjunction with the implementation guide and optimization documents which detail the practical application, including realistic collateral transfer logic.}
\end{abstract}

\section{Introduction}

Perpetual futures contracts use a funding rate mechanism to align prices with the underlying spot asset. This document develops a rigorous mathematical framework for exploiting funding rate inefficiencies, focusing on cross-exchange arbitrage and incorporating advanced modeling techniques.

\section{Perpetual Futures Pricing Model}

Following the random-maturity arbitrage model (He et al., 2024), we formalize the pricing.

\subsection{Notation and Assumptions}

\begin{definition}[Market Setup]
Market includes: Spot price $S_t$, Futures price $F_t$, Risk-free rate $r^f$, Spot borrow/lend rate $r^a$.
\end{definition}

\begin{assumption}[Market Frictions]
Incorporated: Proportional fees $c_s, c_f$; Slippage $s(v,L)$; Hedge holding costs $r^a$; Execution latency $\tau_{exec}$.
\end{assumption}

\subsection{Funding Rate Mechanics}

Payments occur at $t_k = k \Delta T$ based on the average premium/discount over the interval.

\begin{definition}[Funding Rate (Generalized)]
Let $I_t$ be the index/oracle price. Funding rate at $t_{k+1}$ is:
\begin{equation}
    FR_{k+1} = \text{Clamp}_t\left( \text{Avg}_{t \in [t_k, t_{k+1}]} \left( \frac{F_t - I_t}{I_t} \right) \cdot D + R \right)
\end{equation}
where $D$ is dampening factor, $R$ interest component.

\textit{Critique Suggestion:} The clamping function \text{Clamp}_t(\cdot) should ideally be dynamic, adapting to market volatility $\sigma_{premium}$: \text{Clamp Range} = \pm k \cdot \sigma_{\text{premium}}.
\end{definition}

\begin{definition}[Funding Payment]
Payment $FP_{k+1}$ for position size $q$ at $t_{k+1}$:
\begin{equation}
    FP_{k+1} = q \cdot V_{k+1} \cdot FR_{k+1}
\end{equation}
where $V_{k+1}$ is valuation price (mark/oracle).

\textit{Critique Suggestion:} Use multi-oracle average or TWAP for $V_{k+1}$ (if based on oracle) to mitigate manipulation risk.
\end{definition}

\subsection{No-Arbitrage Bounds with Costs}

\begin{theorem}[Refined No-Arbitrage Bounds with Execution Costs]
Cash-and-carry (buy spot, short future), incorporating execution costs (latency $\tau_{exec}$, slippage $s$) and spot costs $r^a$:
\begin{equation}
    (F_t - s_{sell}) e^{-\tau_{exec} r^f} (1-c_f) \le (S_t + s_{buy}) (1+c_s) e^{(r^a - \mathbb{E}_t[\text{Avg}(FR_{pos})])\tau}
\end{equation}
Reverse cash-and-carry (short spot, long future):
\begin{equation}
    (F_t + s_{buy}) e^{-\tau_{exec} r^f} (1+c_f) \ge (S_t - s_{sell}) (1-c_s) e^{(r^a - \mathbb{E}_t[\text{Avg}(FR_{neg})])\tau}
\end{equation}
where $\mathbb{E}_t[FR]$ represents the expected funding rate.

\textit{Critique Suggestion:} Estimating $\mathbb{E}_t[FR]$ accurately is key; consider ML models (e.g., LSTM) for prediction.
\end{theorem}

\section{Single-Exchange Funding Rate Arbitrage}

This strategy involves taking a position in a perpetual future on a single exchange $X$ and hedging the price risk, aiming to capture the funding rate $FR^X$. Assuming a delta-neutral hedge (e.g., holding spot or an offsetting future on another venue with negligible funding), the primary expected profit comes from the funding rate itself, minus costs.

Let $F^X_t$ be the perpetual future price and $S_t$ be the spot price (or hedge instrument price). The basis is $B^X_t = F^X_t - S_t$.
The expected profit $\pi^X_t$ over the next funding period for a position entered at time $t$ is:
\begin{equation} \label{eq:profit_single_strat}
    \mathbb{E}_t[\pi^X_t] \approx \mathbb{E}_t[FR^X_{t+1}] \cdot \text{PositionSize} - C^X_t
\end{equation}
where $\mathbb{E}_t[FR^X_{t+1}]$ is the expected funding rate for the next period, and $C^X_t$ represents the total costs associated with establishing and potentially hedging the position.

$C^X_t$ includes:
\begin{itemize}
    \item Exchange Fees: Taker or maker fees ($c^X$) for entering the perpetual position.
    \item Slippage: Estimated slippage ($s^X$) incurred when entering the position, dependent on size and liquidity. $s^X \approx \beta \cdot (\text{Size} / \text{Liquidity})$.
    \item Hedge Costs: Costs associated with the hedge (e.g., spot trading fees, borrowing costs $r^a$ if shorting spot).
\end{itemize}

\subsection{Kelly Criterion for Single-Exchange Funding Rate}

Assuming the profit $\pi^X_t$ follows a distribution with mean $\mu = \mathbb{E}_t[\pi^X_t]$ and variance $\sigma^2 = \text{Var}_t(\pi^X_t)$, where the variance is primarily driven by the basis volatility $\sigma_{B^X,t}$, the Kelly fraction $f^X_t$ is:
\begin{equation} \label{eq:kelly_basis_singleX_strat}
    f^X_t = \frac{\mathbb{E}_t[\pi^X_t]}{\text{Var}_t(\pi^X_t)} \approx \frac{\mathbb{E}_t[FR^X_{t+1}] \cdot S_t - C^X_t}{(\sigma_{B^X, t})^2}
\end{equation}
where we approximate the variance of the profit by the variance of the basis movement, and normalize the profit by the spot price $S_t$ for the denominator units. $C^X_t$ must also be appropriately scaled relative to the position size represented by the Kelly fraction.

We use a fractional Kelly approach: $f^*_{actual} = \alpha \cdot f^X_t$, where $0 < \alpha \le 1$. \textbf{This fraction is further constrained by available collateral, considering transfer delays, as detailed in FORMULA\_SUMMARY.md and implementation guide: $f_{constrained}$.}

\section{Cross-Exchange Funding Rate Arbitrage}

This strategy involves taking opposing positions on the same perpetual contract (or highly correlated contracts) on two different exchanges, A and B, to capture the difference in their funding rates.

The Net Funding Differential (NFD) between exchange A and exchange B is:
\begin{equation}
    NFD^{AB}_t = FR^A_t - FR^B_t
\end{equation}
The expected profit $\pi^{AB}_t$ over the next funding period is:
\begin{equation} \label{eq:profit_cross_strat}
    \mathbb{E}_t[\pi^{AB}_t] \approx \mathbb{E}_t[NFD^{AB}_{t+1}] \cdot \text{PositionSize} - C^{AB}_t
\end{equation}
where $\mathbb{E}_t[NFD^{AB}_{t+1}]$ is the expected NFD for the next period, and $C^{AB}_t$ represents the total costs of executing the two legs.

$C^{AB}_t$ (Cost-Adjusted NFD components):
\begin{itemize}
    \item Exchange Fees: Sum of fees on both exchanges ($c^A + c^B$).
    \item Slippage: Combined slippage from executing on both exchanges ($s^A + s^B$). $s^{AB} \approx s^A(q, L^A_t) + s^B(q, L^B_t)$.
    \item Latency Cost / Legging Risk ($LC^{AB}$): Implicit cost due to potential adverse price movement between executing leg 1 and leg 2. This can be modeled as an expected value based on short-term volatility and expected execution delay.
    \item Collateral Transfer Costs ($\gamma_{transfer}$): \textbf{Estimated cost of the chosen transfer path (Direct/Bridge) if rebalancing is required for the trade. See collateral optimization logic.}
\end{itemize}

The primary risk is the volatility of the cross-exchange basis, $B^{AB}_t = F^A_t - F^B_t$.

\subsection{Kelly Criterion for Cross-Exchange NFD}

Similar to the single-exchange case, the Kelly fraction $f^{AB}_t$ for the cross-exchange strategy is:
\begin{equation} \label{eq:kelly_cross_pair_strat_updated}
    f^{AB}_t = \frac{\mathbb{E}_t[\pi^{AB}_t]}{\text{Var}_t(\pi^{AB}_t)} \approx \frac{(\mathbb{E}_t[NFD^{AB}_{t+1}] - C^{AB}_{total} / \text{Size}) \cdot S_t}{(\sigma_{B^{AB}, t})^2}
\end{equation}
where $\sigma_{B^{AB}, t}$ is the volatility of the cross-exchange basis $B^{AB}_t$, and \textbf{$C^{AB}_{total}$ now includes estimated transfer costs $\gamma_{transfer}$}. Costs need careful scaling relative to the position size implied by $f^{AB}_t$.

We apply fractional Kelly: $f^*_{actual} = \alpha \cdot f^{AB}_t$. \textbf{This fraction is also subject to the collateral constraint $f_{constrained}$.}

\section{Hyperliquid Funding Rate Mechanisms}

\subsection{Standard Perpetuals}
Index price $I_t$ = Oracle Price (CEX median). Valuation $V_{k+1}$ likely Mark Price.

\subsection{Hyperps (Hyperliquid-only Perps)}
Index price $I^{EMA}_t$ = 8-hour EMA of Hyperliquid mark prices.
\begin{equation}
    FR_{k+1} = \text{Clamp}_t\left( \text{Avg}_{t \in [t_k, t_{k+1}]} \left( \frac{F_t - I^{EMA}_t}{I^{EMA}_t} \right) \cdot D + R \right)
\end{equation}
EMA calculation (corrected interpretation):
\begin{equation}
    I^{EMA}_t = \frac{\sum_{i=0}^{N-1} M_{t-i} w_i}{\sum_{i=0}^{N-1} w_i}, \quad w_i = e^{-i/\tau_{EMA}(t)}
\label{eq:hyperp_ema_adaptive}
\end{equation}
where $\tau_{EMA}(t)$ could be adaptive based on volatility $\sigma_t$. Safeguard cap $\le M_{initial} \times 4$ might be replaced by dynamic threshold (e.g., $\pm 3\sigma$). Dampening factor $D$ might be smaller.

\section{Modeling Considerations & Advanced Concepts}

\begin{itemize}
    \item \textbf{Stochastic/Adaptive Models:} Use OU, HMM, or ML (e.g., LSTM) for adaptive $\mathbb{E}_t[FR]$. Use GARCH/Stochastic Vol models for adaptive $\sigma_{B,t}$.
    \item \textbf{Optimal Dynamic Hedging:} Employ time-varying $\delta_t^*$.
    \item \textbf{Execution Costs:} Explicitly model variable fees, slippage $s(q,L)$, and latency costs in profit/NFD.
    \item \textbf{Risk Management:} Use CVaR alongside VaR. Consider Sortino/Omega ratios. Base position limits on portfolio risk ($\sqrt{w^T \Sigma w}$) and use fractional, \textbf{transfer-constrained} Kelly.
    \item \textbf{Parameter Uncertainty:} Use Bayesian estimation, shrinkage, robust optimization.
    \item \textbf{Collateral Management:} Implement \textbf{dynamic path selection, bridge integration, caching, ML targets, retries, contingency borrowing} as detailed in implementation guide and diagrams.
\end{itemize}

\section{Conclusion}

This enhanced framework incorporates critique suggestions: adaptive parameters (funding rate expectations, volatility, EMA decay, clamping), explicit execution costs (\textbf{including transfers}), advanced risk metrics (CVaR), and robust estimation notes. \textbf{The Kelly criterion is refined to include transfer costs and collateral constraints.} These refinements aim for a more realistic and robust representation of funding rate arbitrage, especially for dynamic crypto markets. \textbf{Practical implementation relies heavily on the detailed collateral management and execution logic outlined in accompanying documents.}

\end{document}